% !TEX root = ../main.tex
\chapter{Metodología y Puntos a tratar}
\label{cha:MetodologíayPaT}

En esta sección definiremos la metodología que realizaremos y los pasos que seguiremos para realizar este trabajo como la selección del modelo que utilizaremos, las herramientas, los datos entre otros aspectos relevantes para el desarrollo del trabajo de investigación, pasos correspondientes al trabajo necesario para cumplir con los objetivos específicos y que nos permitirán llegar al objetivo general del proyecto.

\begin{enumerate}
    \item Investigación de estado del arte y revisión bibliográfica de conceptos que se tratarán.
    \item Definición del enfoque de investigación.
    \item Descripción del dataset a utilizar.
    \item Investigación de modelos NLU (Monolingües y Multilingües) y selección de los modelos a utilizar.
    \item Definición de las herramientas y técnicas a utilizar.
    \item Preparación de ambientes de desarrollo 
    \item Preparación y preprocesamiento de datos
    \item Verificación de la calidad de los datos preprocesados
    \item Procesamiento de datos con modelo NLU seleccionado
    \item Análisis de datos obtenidos
    \item Descripción de trabajo futuro
\end{enumerate}

\section{\textbf{Investigación de estado del arte:}}
Primero se comenzó por ingvestigar a fondo las temáticas a tratar, herramientas, contexto, conceptos que luego irian en el marco teórico, modelos NLU, técnicas de preprocesamiento de datos, análisis de datos, entre otros aspectos ecesarios para el desarrollo de la investigación.

\section{\textbf{Definición del enfoque de investigación:}}
Se definió si el enfoque iba a ser cualitativo, cuantitativo o mixto, considerando cosas como las características y/o columnas del dataset, los objetivos planteados y las técnicas a utilizar, para poder desarrollar un fundamento mas sólido.

\section{\textbf{Descripción del dataset a utilizar:}}
Se estudia y describe la estructura del dataset a utilizar, columnas, caracteristicas y todos los datos necesarios para comprender los datos con los que se están trabajando y cuales realmente representan un aporte al trabajo y cuales no.


\section{\textbf{Investigación de modelos NLU:}}
Sección en la que se estudia un pequeño estado del arte de los modelos NLU, origen arquitectura, linea temporal y para que sirven o con que datos están entrenados,

\section{\textbf{Definición de las herramientas y técnicas a utilizar:}}Sección en la que se fedinien las diferenes hjerramientas que utilizaraán para preprocesar los datos, analizar el lenguaje natural y realizar el análisis de datos posterior.


\section{\textbf{Preparación de ambientes de desarrollo:}}
Se preparan dependencioas, embnientes, repositorios, y los espacios necesarios oara la realización del trabajo.

\section{\textbf{Preparación y preprocesamiento de datos}}
Etapa en la que realizaremos la limpieza, etiquetado, aseguramiento del mismo y preparación de los datos para que estos puedan ser analizados por el modelo NLU seleccionado posteriormente, además de reducir la cantidad de ruido en los datos y mejorar la calidad de los mismos para obtener mejores resultados (técnica que ya fue mencionada anteriormente).\\
Además es la sección que determina con los datos finales con los que se deberían trabajar, siendo esta una de las etapas mas significativas e imprescindibles del proyecto, donde se comprenden etapas como el reconocimiento de lenguaje, limpieza de datos, etiquetado o corroboración manual. Resultando todo esto en los datos etiquetados para ya pasar por un modelo mas robusto.

\section{\textbf{Verificación de la calidad de los datos preprocesados}}
Una vez realizado el proceso de preprocesado de los datos se debe valildar que estos mismos posean cierto grado de fidelidad y que no sean datos aparentemente correctos pero que en realidad no lo sean. Se planteó la idea de revisiones con técnicas como Embeddings con comparación por similitud del coseno o modelos LLM como Groq, modelos que de por si resultan mas robustos y fueron preentrenados con un gran corpus de datos, de manera que (pese a su coste computacional) nos serían posiblemente muy útiles para validar la calidad de los datos preprocesados/etiquetados.

\section{\textbf{Procesamiento de datos con modelo NLU seleccionado}}

Habiendo ya preprocesados los datos y determinado que su calidad es la adecuada para poder continuar con la operación, podemos proceder con la utilización de un modelo NLU multilingue que realmente pueda etiquetar los datos correctamente (no trabajando de manera excluyente a lo anterior, si no que complementandolo), siendo este el último y gran filtro para determinar si un comentario es o no una sugerencia operativa, y así finalmente despues de suficientes controles de calidad, obtener los datos finales que nos permitirán realizar el análisis estadístico y comparativo necesario para llegar a conclusiones relevantes.

\section{\textbf{Análisis de datos obtenidos}}

Una vez ya realizadas las pruebas y obtenidos los datos, se deberá realizar un análisis estadístico del rendimiento general, tablas de pruebas donde se demuestre las distintas comparaciones de eficacia de ya sea modelos NLU (considerados a parte si es que aplica), técnicas de preprocesamiento frente a otros trabajos de investigación y las lecciones aprendidas en el ejercicio del trabajo realizado.

\section{\textbf{Descripción de trabajo futuro}}

Luego de finiquitado el trabajo como tal y haber dado las respectivas conclusiones respecto al trabajo realizado, se deberá plantear una sección donde se mencione mínimamente lo relacionado a la manera en que se trabajó en la investigación, pero principalmente se deberá plantear el trabajo futuro que podría realizarse a partir de lo ya realizado, otras formas de mejorar la detección de las sugerencias operacionales, nuevas técnicas, nuevos modelos NLU, LLMs, etc.